\chapter{Cronograma Previsto}
\label{cap:cronograma}

A revisão bibliográfica e o desenvolvimento do protótipo encontram-se concluídos. O ambiente virtual foi implementado --- com os sistemas principais, a interface, a lógica de pontuação, o monitoramento de desempenho e o registro de eventos --- e já passou pela validação técnica descrita no Capítulo~\ref{cap:validacao}. Restam ajustes de experiência do usuário e otimização, sem alteração do fluxo de interação, já finalizado.

O tempo restante concentra-se na avaliação pedagógica com participantes, cujo caminho crítico depende de fatores externos. A defesa final tem prazo de titulação rígido em novembro de 2026. A Tabela~\ref{tab:cronograma} distribui as etapas pendentes entre junho e novembro de 2026.

\begin{table}[htb]
\centering
\caption{Cronograma das etapas restantes (2026).}
\label{tab:cronograma}
\begin{tabular}{|l|c|c|c|c|c|c|}
\hline
\textbf{Etapa} & \textbf{Jun} & \textbf{Jul} & \textbf{Ago} & \textbf{Set} & \textbf{Out} & \textbf{Nov} \\
\hline
Validação do teste de conhecimento & X & X & & & & \\
\hline
Ajustes de UX e otimização & X & X & & & & \\
\hline
Aprovação no CEP (instrumento anexado) & & X & & & & \\
\hline
Recrutamento e sessões (Camada 1) & & & X & & & \\
\hline
Teste de retenção (condicional) & & & & X & X & \\
\hline
Análise de dados & & & & & X & X \\
\hline
Redação e defesa & & & & & & X \\
\hline
\end{tabular}
\end{table}

O item de maior risco é a validação de conteúdo do teste de conhecimento, que também serve de base ao teste de retenção. Por isso, ela é prioridade imediata e corre em paralelo aos ajustes de interface. O instrumento será anexado ao protocolo do Comitê de Ética em Pesquisa antes de sua aprovação, evitando um ciclo adicional de emenda. Os questionários SUS e IPQ já estão prontos.

A coleta organiza-se em camadas, de modo que a defesa não dependa da etapa de maior incerteza de calendário. A camada primária --- ganho de conhecimento (pré e pós-teste), motivação (RIMMS), usabilidade (SUS) e presença (IPQ) --- é coletada em sessão única e ancora a avaliação. O teste de retenção, aplicado de 30 a 45 dias após o treinamento, é condicional: se não couber no prazo rígido de novembro, será declarado como limitação e recomendado para trabalho futuro, sem comprometer as demais análises.
